\section{Relaci\'on t\'ermica intr\'inseca de la escala \(108\)}

La estructura cronol\'ogica \(108\) fija

\begin{equation}
k_B\Theta_{\rm clk}\log3
=\frac{h_{\rm int}}{108t_0}
=\frac{2\pi\hbar_{\rm int}}{108t_0}.
\label{eq:boltzmann-clock}
\end{equation}

Esta igualdad no determina por separado una coordenada num\'erica de \(k_B\) y otra de \(\Theta_{\rm clk}\) sin una normalizaci\'on adicional. S\'i determina su producto como secci\'on de energ\'ia. El factor \(\log3\) tiene tres antecedentes compatibles:

\begin{equation}
\log3=C_1+C_2-2C_0
=\frac{S_{\TRIT}}{k_B}
=\frac{E_P}{k_B\Theta_{\rm clk}}.
\label{eq:log3-triple}
\end{equation}

Las tres igualdades corresponden a realizaciones residual, informacional y
termodin\'amica de un mismo escalar. La coincidencia del valor no identifica
sus dominios respectivos.

\section{Formulaciones espectrales de la ley de desplazamiento de Wien}

Para la densidad espectral por longitud de onda, el m\'aximo satisface

\[
5(1-\ee^{-x_5})=x_5,
\qquad
x_5=5+W_0(-5\ee^{-5}).
\]

Con \cref{eq:boltzmann-clock},

\begin{equation}
\lambda_{\max}(\Theta_{\rm clk})
=\frac{108\log3}{x_5}\,\ell_0.
\label{eq:wien-lambda}
\end{equation}

Para la densidad por frecuencia,

\[
3(1-\ee^{-x_3})=x_3,
\qquad
x_3=3+W_0(-3\ee^{-3}),
\]

y

\begin{equation}
\nu_{\max}(\Theta_{\rm clk})
=\frac{x_3}{108t_0\log3}.
\label{eq:wien-frequency}
\end{equation}

Los m\'aximos no son rec\'iprocos, puesto que las densidades espectrales se
transforman con el jacobiano correspondiente. El cociente \(x_5/x_3\)
pertenece a la estructura de la parametrizaci\'on y no representa una
inconsistencia.

\section{Ley de Stefan--Boltzmann y termodin\'amica del gas de fotones}

Evaluar la ley de radiaci\'on en \(\Theta_{\rm clk}\) produce el flujo intr\'inseco

\begin{equation}
j_\star(\Theta_{\rm clk})
=\frac{2\pi^5}{15\,108^4(\log3)^4}\,
\frac{h_{\rm int}}{\ell_0^2t_0^2},
\label{eq:stefan-hmt}
\end{equation}

en la normalizaci\'on declarada. La constante de Stefan--Boltzmann no se inserta como primitiva: resulta de integrar el espectro bos\'onico una vez fijados acci\'on, propagaci\'on y temperatura.
Las leyes espectrales empleadas son las de Planck y Wien \cite{Planck,Wien}; HMT fija aqu\'i la secci\'on sobre la que se eval\'uan, no sus constantes por contraste retrospectivo.

En unidades de la escala de Planck,

\begin{align}
n_\gamma\ell_P^3
&=\frac{2\zeta(3)}{\pi^2\log^3 3},
\label{eq:photon-number}\\
\frac{u_\gamma\ell_P^3}{E_P}
&=\frac{\pi^2}{15\log^4 3},
\label{eq:photon-energy}\\
\frac{s_\gamma\ell_P^3}{k_B}
&=\frac{4\pi^2}{45\log^3 3}.
\label{eq:photon-entropy}
\end{align}

\(\zeta(3)\) representa el momento bos\'onico de orden tres, mientras que
\(\log3\) determina la energ\'ia de decisi\'on TRIT. La ecuaci\'on conserva la
distinci\'on entre ambos invariantes.

La conductancia t\'ermica cu\'antica se expresa como

\begin{equation}
\frac{\kappa_Qt_0}{k_B}=\frac{\pi^2}{324\log3}.
\label{eq:thermal-conductance}
\end{equation}

\section{Confluencia metrológica de los invariantes térmicos y radiativos}

\begin{longtable}{@{}p{0.18\textwidth}p{0.27\textwidth}p{0.45\textwidth}@{}}
\toprule
Objeto & Ecuaci\'on interna & Significado estructural\\
\midrule
\endhead
\(R_K\) & \(Z_{\rm vac}/(2\alpha)\) & resistencia cu\'antica bidireccional invariante\\
\(G_0\) & \(4\alpha/Z_{\rm vac}\) & conductancia complementaria\\
\(\Phi_0\) & \(\sqrt{hZ/(8\alpha)}\) & flujo que conserva orientaci\'on de acci\'on\\
\(K_J\) & \(\Phi_0^{-1}\) & relaci\'on rec\'iproca entre frecuencia y flujo\\
\(E_P\) & \(k_B\Theta\log3\) & energ\'ia del ciclo y contenido informacional TRIT\\
Wien & \(108\log3/x_5\) & extremo espectral en la escala cronol\'ogica\\
Ap\'ery fot\'onica & \(2\zeta(3)/(\pi^2\log^3 3)\) & densidad de n\'umero bos\'onica\\
\bottomrule
\end{longtable}

\begin{remark}[Control negativo]
El bloque queda refutado si un cambio de variable espectral conserva
indebidamente el mismo extremo de Wien, si \(\zeta(3)\) sustituye a
\(\log3\), o si se asignan valores separados a \(k_B\) y \(\Theta\) sin
especificar una carta. Las identidades
\cref{eq:quantum-electrical-identities,eq:planck-trit} proporcionan controles
algebraicos directos.
\end{remark}

\subsection*{Alcance lógico y dominio de validez}
Cuantos el\'ectricos y confluencia Planck--TRIT: \emph{Teorema interno exacto}. Wien, Stefan--Boltzmann y gas fot\'onico: \emph{Composici\'on matem\'atica exacta} con el continuo bos\'onico declarado. La carta SI permanece como \emph{Contraste metrol\'ogico externo}.

